\begin{table}[!htbp]
  \centering
  \caption{不同量化路径的离线产物、运行时缓存载荷与默认执行后端}
  \label{tab:ch3-runtime-paths}
  \scriptsize
  \setlength{\tabcolsep}{2.0pt}
  \renewcommand{\arraystretch}{1.14}
  \begin{tabularx}{\linewidth}{>{\raggedright\arraybackslash}p{1.35cm}>{\raggedright\arraybackslash}p{2.35cm}>{\raggedright\arraybackslash}p{2.95cm}>{\raggedright\arraybackslash}p{2.25cm}X}
    \toprule
    路径 & 离线产物 & 追加写入 & 解码后端 & 边界说明 \\
    \midrule
    INT8 & scale、保护裕度、轴元数据 & INT8 K/V 与 scale & Triton 融合后端 & 主部署路径 \\
    对称 INT4 & scale、轴元数据 & packed INT4 K/V & 参考实现或 Triton 试探 & 压力测试；解包计入成本 \\
    RoleAlign & $(p_K,p_V)$、K/V 轴、头数 & K per-channel；V per-token & 参考实现；融合扩展 & 质量核对优先；速度见第四章 \\
    \bottomrule
  \end{tabularx}
\end{table}
